\documentclass[11pt]{article}
\usepackage{amsmath,amssymb,booktabs}
\usepackage[margin=1in]{geometry}
\title{Step 202: Numerical Attempt on \(\Xi_{\mathrm{matrix\_source}}\)}
\date{2026-05-16}
\begin{document}
\maketitle

\section*{Verdict}

\[
\boxed{\texttt{V\_xi\_matrix\_source\_partial}.}
\]

The Step 201 Burnol unblock is used.  The numerical pipeline evaluates
\(E_{1/2}\) and assembles the literal de Branges kernel matrix
\[
G_{ij}=K_{1/2}^{\Gamma}(\rho_j,\rho_i).
\]
The residual \(\Xi_{\mathrm{matrix\_source}}\) is not decided, because the
literal convention gives zero diagonal entries \(G_{ii}=0\) on the critical
line, so the normalized vectors \(e_i=\kappa_i/\sqrt{G_{ii}}\) cannot be
formed without the correct diagonal reproducing-kernel limit or
inner-product convention.

\section*{Inherited Paper-Grounded Formulas}

Burnol 2002 (math/0208121), Theorem 4, gives the explicit projection
\(\pi_\lambda=P_{K_\lambda}=P_{L_a^\Gamma}\).  Corollary 5 gives the
\((1\pm F_\lambda)^{-1}\) simplifications.  Definition 2 gives
\(\psi_\pm^\lambda\).  Theorem 8 gives
\[
E_\lambda(w)=\pi^{-w/2}\Gamma(w/2)\left[
\lambda^{1/2-w}+\frac{\sqrt{\lambda}}{2}
\int_\lambda^\infty
(\psi_+^\lambda(t)-\psi_-^\lambda(t))t^{-w}\,dt\right],
\qquad \Re w>0.
\]
Burnol 2002, equation 1, gives
\[
K_{B(E_\lambda)}(z_1,z_2)=
\frac{E_\lambda(z_1)\overline{E_\lambda(z_2)}
-E_\lambda(1-z_1)\overline{E_\lambda(1-z_2)}}
{z_1+z_2-1}.
\]
Burnol 2004 (math/0112254), Definition 6.1 and Theorem 6.10, supply the
carrier \(K_\lambda\) and the completion
\(\pi^{-s/2}\Gamma(s/2)\).

\section*{Numerical Method}

For \(\lambda=1/2\), the script discretizes the cosine-transform operator on
\([0,1/2]\), solves
\[
(1+F_{1/2})u_+=2\cos(\pi y),\qquad
(1-F_{1/2})u_-=2\cos(\pi y),
\]
and evaluates
\[
\psi_+^{1/2}(t)-\psi_-^{1/2}(t)
=-F_+u_+(t)-F_+u_-(t)
\]
inside Burnol's Theorem 8 integral.  Settings were \(70\) scalar decimal
digits, \(160\) cosine-resolvent nodes, \(520\) Gauss--Legendre integral
nodes, cutoff \(80\), and check cutoff \(55\).

\section*{Computed \(E_{1/2}\) Values}

\begin{center}
\begin{tabular}{c c c}
\toprule
point & \(E_{1/2}(w)\) & error bound\\
\midrule
\(\rho_1\) & \(-4.08121209197\cdot 10^{-6}+1.57450659741\cdot 10^{-5}i\) & \(1.66\cdot 10^{-8}\)\\
\(\rho_2\) & \(-5.56494927209\cdot 10^{-8}-5.32637080111\cdot 10^{-8}i\) & \(1.11\cdot 10^{-9}\)\\
\(\rho_3\) & \(-2.85254026901\cdot 10^{-9}+3.56620933992\cdot 10^{-10}i\) & \(1.06\cdot 10^{-9}\)\\
\bottomrule
\end{tabular}
\end{center}

The values at \(1-\rho_i\) are the conjugate values in this discretization.

\section*{Gram Matrix Layer}

The literal \(G\)-matrix has zero diagonal entries:
\[
|G|\approx
\begin{pmatrix}
0 & 6.2212\cdot10^{-14} & 2.2203\cdot10^{-15}\\
6.2212\cdot10^{-14} & 0 & 7.4635\cdot10^{-18}\\
2.2203\cdot10^{-15} & 7.4635\cdot10^{-18} & 0
\end{pmatrix}.
\]
This is not a usable Hilbert Gram matrix for normalizing the \(e_i\).  The
matrix-source residual therefore exposes a sharper sub-residual:
determine the correct critical-line kernel convention and diagonal limit for
the zero-evaluator atoms.

\section*{Commutator Matrix}

The desired entries are
\[
c_{ij}(\ell)=
\left\langle e_i,(I-P_\infty)M_{m_\ell}P_\infty e_j\right\rangle,
\qquad
m_\ell(s)=\exp[-\ell(1/2-s)].
\]
They remain explicit but uncertified numerically.  A certified evaluation must
first resolve \(e_i\), then compute transported boundary samples
\(\kappa_i=T_{1/2}^*K_{1/2}^\Gamma(\cdot,\rho_i)\) and include the full
\(K_\infty^{op}\) PSWF correction.

\section*{Nonclaim}

This step does not prove RH, does not close \(\Xi_{BC}\), and does not close
Branch B.  It identifies the next numerical/theoretical blocker after the
Step 201 Burnol \(\kappa\) unblock.

\end{document}

